% SOURCE: docs/比赛材料/开发目的.md + docs/比赛材料/设计原理与方法.md §引言
\section{选题意义与目标定位}

\subsection{混沌理论的教学价值与现状}

非线性动力学与混沌理论是 20 世纪物理学的三大革命性突破之一，与相对论和量子力学并列。
混沌揭示了确定性系统产生表观随机行为的深刻机制——这一发现从根本上改变了人类对“可预测性”的哲学理解。
然而，在国内本科物理课程体系中，混沌内容长期处于边缘地位：多数《大学物理》课程以线性振动（简谐运动、阻尼振动、受迫振动）为核心，
对非线性振动和混沌仅以“科普讲座”或“绪论末尾”形式简要提及，缺乏系统性教学支撑。

学生的典型认知障碍包括：(1) 将“决定论”与“可预测性”等同——认为既然微分方程唯一确定，则长期预测必为可行；
(2) 将“偶然性”与“不可预测性”混为一谈——无法区分量子力学的不确定性与经典混沌的初值敏感性。
这些概念混淆是混沌教学中的主要卡点，传统以文字描述和静态图片为主的教材呈现方式难以有效纠正。

\subsection{实体实验的固有瓶颈与数字资源的不足}

双摆是混沌教学的经典模型——仅两个自由度却能展现从周期运动到完全混沌的全谱系动力学行为。
但实体双摆实验面临多重困境：

\begin{enumerate}[label=(\arabic*), leftmargin=*]
  \item \textbf{加工精度要求苛刻}：转轴摩擦、空气阻力、轴承间隙会掩盖真实的混沌动力学行为，使实验现象偏离理论预期；
  \item \textbf{参数调节效率低下}：摆长、质量需物理拆装，无法实现“滑杆拖拽即生效”的即时反馈，而绘制一条完整分岔图需扫描数百个参数点，课堂时间内无法完成；
  \item \textbf{数据采集存在系统性偏差}：传感器精度、量化误差、机械摩擦损耗使测量结果与理论模型之间存在不可忽略的差距；
  \item \textbf{初始条件难以精确复现}：实验可重复性差，参数调试不当会掩盖混沌态中的周期窗口、拉长实验周期。
\end{enumerate}

另一方面，现有数字资源也未能填补空白：Flash/MP4 演示动画“只能看不能操作”，MATLAB/Python 离线脚本“需安装环境、调试代码，门槛高”。
多数仿真工具停留在现象演示层面，缺乏从“感知→分析→验证→创造”的完整认知路径，且商业软件授权复杂、成本高，不适合大规模课堂部署。

\subsection{本作品的目标定位}

本作品旨在提供一个\textbf{纯浏览器端、零安装门槛、离线可运行}的双摆非线性动力学虚拟仿真平台，
通过高精度数值积分引擎驱动实时物理仿真，以四层递进认知架构覆盖从入门感知到科研训练的完整教学需求。

具体目标包括：
\begin{enumerate}[label=(\arabic*), leftmargin=*]
  \item \textbf{以数值仿真替代实体实验}：自实现 RKF45 自适应步长积分引擎，在 Web Worker 中并行求解双摆 ODE，消除机械摩擦、传感器噪声和量化误差；
  \item \textbf{构建完整认知闭环}：设计探索（感知）、分析（理解）、实验（验证）、故事（展示）四种工作模式，使学生从直观感受混沌现象出发，逐步建立定量分析能力；
  \item \textbf{实现零门槛部署}：纯客户端架构，产出物为自包含 HTML5 文件夹，支持 \texttt{file://} 协议离线运行，可嵌入课程网站、拷入 U 盘、部署在教室一体机或科技馆终端；
  \item \textbf{纠正关键认知误区}：通过蝴蝶效应双视口对比器（$\delta\theta_1 = 10^{-6\circ}$）和时间反演双模式实验（精确/数值对比），直观展示“确定性但不可预测”的物理本质。
\end{enumerate}

\endinput
